% This file is included by uet_thesis.tex.
% Replace blue placeholders and remove example text before submission.

\chapter{Planning, Risk, and Traceability}
\label{chap:planning}
% =============================================================================

\section{Design Traceability Matrix}
\begin{guidance}
Demonstrate that objectives are supported by approved scope, requirements, design,
and planned verification. Add rows as needed; a requirement may appear in more than
one row. Use exact identifiers and reference detailed tests in the separate testing
document rather than copying them here.
\end{guidance}

{\footnotesize
\begin{longtable}{@{}p{0.12\textwidth}p{0.11\textwidth}p{0.16\textwidth}p{0.23\textwidth}p{0.24\textwidth}@{}}
\caption{Design traceability matrix}\label{tab:traceability}\\
\toprule
\textbf{Objective} & \textbf{Scope} & \textbf{Requirement} & \textbf{Design element} & \textbf{Testing-document reference} \\
\midrule
\endfirsthead
\toprule
\textbf{Objective} & \textbf{Scope} & \textbf{Requirement} & \textbf{Design element} & \textbf{Testing-document reference} \\
\midrule
\endhead
OBJ-01 & SCP-01 & FR-01 & UC-01; appointment sequence; Appointment class & \placeholder{TEST-FR-01 / Section X} \\
\placeholder{OBJ-XX} & \placeholder{SCP-XX} & \placeholder{FR-XX} & \placeholder{Diagram, component, model, or decision ID} & \placeholder{Test or validation reference} \\
\bottomrule
\end{longtable}
}

\section{Verification and Validation Strategy}
\begin{guidance}
Summarize how the team will verify that the system conforms to its specifications
and validate that it addresses stakeholder needs or research objectives. Identify
the major methods and responsible roles only; place detailed procedures, cases,
data, expected results, and actual results in the testing document.
\end{guidance}

\begin{table}[H]
\centering
\caption{Verification and validation strategy}
\label{tab:vv-strategy}
\begin{tabularx}{\textwidth}{@{}p{0.20\textwidth}X X p{0.20\textwidth}@{}}
\toprule
\textbf{Target} & \textbf{Method} & \textbf{Evidence} & \textbf{Responsible role} \\
\midrule
Functional requirements & Functional and integration testing & Approved test results in the testing document & \placeholder{Team role} \\
Stakeholder needs & \placeholder{Demonstration, usability study, review, or field evaluation} & \placeholder{Feedback or validation record} & \placeholder{Team role/stakeholder} \\
Research findings & \placeholder{Statistical validation, cross-validation, triangulation, or expert review} & \placeholder{Analysis outputs} & \placeholder{Responsible role} \\
\bottomrule
\end{tabularx}
\end{table}

\section{Testing-Document Reference}
\begin{guidance}
Identify the separate testing document that contains the detailed acceptance
criteria, test cases, datasets, procedures, results, and defect records. Keep its
version consistent with this design document.
\end{guidance}

\begin{table}[H]
\centering
\caption{Software testing document reference}
\label{tab:testing-reference}
\begin{tabularx}{\textwidth}{@{}p{0.25\textwidth}X@{}}
\toprule
\textbf{Item} & \textbf{Reference} \\
\midrule
Document title & \placeholder{Project Software Testing Document} \\
Version and date & \placeholder{Version; Day Month Year} \\
Repository or submission location & \placeholder{Controlled path, repository location, or submission reference} \\
Responsible team member(s) & \placeholder{Name(s) and role(s)} \\
\bottomrule
\end{tabularx}
\end{table}

\section{Risk Assessment}
\begin{guidance}
Record significant technical, research, ethical, schedule, resource, and external
risks. Rate probability and impact using a consistent scale, assign an owner, and
state a practical response.
\end{guidance}

\begin{table}[H]
\centering
\caption{Risk register}
\label{tab:risk-register}
\small
\begin{tabularx}{\textwidth}{@{}p{0.11\textwidth}X p{0.14\textwidth}p{0.11\textwidth}X p{0.14\textwidth}@{}}
\toprule
\textbf{ID} & \textbf{Risk} & \textbf{Likelihood} & \textbf{Impact} & \textbf{Response} & \textbf{Owner} \\
\midrule
RSK-01 & External API becomes unavailable. & Medium & High & Use an interface abstraction and maintain a documented fallback dataset or mock service. & \placeholder{Role} \\
RSK-02 & \placeholder{Risk event} & \placeholder{Level} & \placeholder{Level} & \placeholder{Avoid, reduce, transfer, or accept} & \placeholder{Role} \\
\bottomrule
\end{tabularx}
\end{table}

\section{Project Management Approach and Schedule}
\begin{guidance}
State the development or research management approach and show major phases,
dependencies, milestones, review points, and deliverables. A Gantt chart may be
added when required; keep detailed weekly task tracking in the selected management
tool rather than duplicating it here.
\end{guidance}

\begin{table}[H]
\centering
\caption{Milestones and schedule}
\label{tab:schedule}
\begin{tabularx}{\textwidth}{@{}p{0.14\textwidth}X p{0.20\textwidth}p{0.19\textwidth}p{0.16\textwidth}@{}}
\toprule
\textbf{ID} & \textbf{Milestone or deliverable} & \textbf{Target date} & \textbf{Dependency} & \textbf{Owner} \\
\midrule
MS-01 & \placeholder{Approved requirements baseline} & \placeholder{Date} & \placeholder{None or ID} & \placeholder{Role} \\
MS-02 & \placeholder{Prototype, study, evaluation, or release milestone} & \placeholder{Date} & \placeholder{MS-01} & \placeholder{Role} \\
\bottomrule
\end{tabularx}
\end{table}

\section{Team Roles and Responsibilities}
\begin{guidance}
Assign clear ownership while recognizing shared review and integration work. Roles
should reflect actual responsibilities and not be used to divide the project into
unconnected individual submissions.
\end{guidance}

\begin{table}[H]
\centering
\caption{Team roles and responsibilities}
\label{tab:team-roles}
\begin{tabularx}{\textwidth}{@{}p{0.24\textwidth}p{0.24\textwidth}X@{}}
\toprule
\textbf{Member} & \textbf{Role} & \textbf{Principal responsibilities} \\
\midrule
\placeholder{Name and roll number} & \placeholder{Role} & \placeholder{Owned tasks, reviews, and deliverables} \\
\bottomrule
\end{tabularx}
\end{table}

\section{Source-Code, Data, and Artifact Management}
\begin{guidance}
State how code, documents, datasets, models, diagrams, builds, and other artifacts
will be named, versioned, reviewed, backed up, licensed, and accessed. Identify how
sensitive data and credentials will be kept out of the source repository.
\end{guidance}

\begin{table}[H]
\centering
\caption{Artifact management}
\label{tab:artifact-management}
\begin{tabularx}{\textwidth}{@{}p{0.21\textwidth}X X X@{}}
\toprule
\textbf{Artifact} & \textbf{Storage and versioning} & \textbf{Access or protection} & \textbf{Backup or retention} \\
\midrule
Source code & \placeholder{Repository and branching/review approach} & \placeholder{Access control and secret handling} & \placeholder{Backup or release policy} \\
\placeholder{Data, models, and documents} & \placeholder{Location and versioning} & \placeholder{Access rule} & \placeholder{Retention rule} \\
\bottomrule
\end{tabularx}
\end{table}

\section{Tools and Technologies}
\label{sec:tools}
\begin{guidance}
List only tools and technologies that are planned or already used. State each
item's purpose and relevant version. Explain a selection decision in
Section~\ref{chap:design} only when alternatives and trade-offs matter.
\end{guidance}

\begin{table}[H]
\centering
\caption{Tools and technologies}
\label{tab:tools-technologies}
\begin{tabularx}{\textwidth}{@{}p{0.24\textwidth}p{0.20\textwidth}X X@{}}
\toprule
\textbf{Tool or technology} & \textbf{Version} & \textbf{Purpose} & \textbf{Project use} \\
\midrule
\placeholder{Framework, engine, language, library, IDE, modeling tool, or platform} & \placeholder{Version} & \placeholder{General purpose} & \placeholder{Specific use in this project} \\
\bottomrule
\end{tabularx}
\end{table}
